\documentclass[a4paper,12pt]{article}
\usepackage{color}
\usepackage{graphicx, gensymb}
\usepackage{amsfonts, cancel, amssymb}
\usepackage{amsmath, amsthm, array, esvect, esint}
\usepackage{typearea}
\usepackage[landscape=true]{geometry}
\newcommand\numberthis{\addtocounter{equation}{1}\tag{\theequation}}
\newcommand{\bra}{}
\newcommand{\kom}{}
\newcommand{\brak}{}
\newcommand{\braket}{}
\newcommand{\intx}{}
\newcommand{\intr}{}
\newcommand{\kla}{}
\newcommand{\klam}{}
\newcommand{\klammer}{}
\newcommand{\oper}{}
\newcommand{\tecxt}{}
\newcommand{\Bra}{}
\newcommand{\Ket}{}

\begin{document}

\title{07 Schwarzschild-Metrik}
\author{Joachim Schwardt}
\date{\today}
\maketitle


\section*{Aufgabe 16: Lichablenkung}
%https://aapt.scitation.org/doi/10.1119/1.4866274
Die Schwarzschild-Metrik lautet
\begin{align*}
\tecxt\text{d}s^2&=\kla\left( {1-\frac{R_S}{r}} \right)\tecxt\text{d}t^2-\kla\left( {1-\frac{R_S}{r}} \right)^{-1}\tecxt\text{d}r^2-r^2\tecxt\text{d}\Omega^2
\end{align*}
mit dem Schwarzschild-Radius $R_S=\frac{2G_NM}{c^2}$. Die Geodäten-Gleichung der $\theta$-Komponente wird (unter anderem) durch $\theta=\frac{\pi}{2}=\tecxt\text{const}$ gelöst, die Bewegung verläuft also in einer Ebene:
\begin{align*}
0&=\frac{\tecxt\text{d}^2\theta}{\tecxt\text{d}s^2}+\Gamma^\theta_{\mu\nu}\frac{\tecxt\text{d}x^\mu}{\tecxt\text{d}s}\frac{\tecxt\text{d}x^\nu}{\tecxt\text{d}s}=\frac{\tecxt\text{d}^2\theta}{\tecxt\text{d}s^2} + \frac{2}{r}\frac{\tecxt\text{d}r}{\tecxt\text{d}s}\frac{\tecxt\text{d}\theta}{\tecxt\text{d}s} - \sin\theta\cos\theta\kla\left( {\frac{\tecxt\text{d}\phi}{\tecxt\text{d}s}} \right)^2
\end{align*}
Dabei wurden die Christoffel-Symbole $\Gamma^\theta_{r\theta}=\Gamma^\theta_{\theta r}=\frac{1}{r}$ und $\Gamma^\theta_{\phi\phi}=-\sin\theta\cos\theta$ verwendet. Die Metrik vereinfacht sich dann zu $\tecxt\text{d}s^2=\dots -r^2\tecxt\text{d}\phi^2$.\\ 
Für masselose Teilchen gilt $\tecxt\text{d}s^2=0$, die Bewegungsgleichung wird stattdessen durch einen beliebigen Parameter $\lambda$ charakterisiert. Die Lagrange-Funktion lautet dann 
\begin{align*}
L&=g_{\mu\nu}\frac{\tecxt\text{d}x^\mu}{\tecxt\text{d}\lambda}\frac{\tecxt\text{d}x^\nu}{\tecxt\text{d}\lambda}=\kla\left( {1-\frac{R_S}{r}} \right)\kla\left( {\frac{\tecxt\text{d}t}{\tecxt\text{d}\lambda}} \right)^2-\kla\left( {1-\frac{R_S}{r}} \right)^{-1}\kla\left( {\frac{\tecxt\text{d}r}{\tecxt\text{d}\lambda}} \right)^2-r^2\kla\left( {\frac{\tecxt\text{d}\phi}{\tecxt\text{d}\lambda}} \right)^2
\end{align*}
und hängt nicht explizit von $t$ oder $\phi$ ab. Also gibt es zwei Erhaltungsgrößen:
\begin{align*}
\frac{\tecxt\text{d}}{\tecxt\text{d}\lambda}\frac{\partial L}{\partial \frac{\tecxt\text{d}t}{\tecxt\text{d}\lambda}}&=\frac{\partial L}{\partial t}=0&\Rightarrow &&\frac{\partial L}{\partial \frac{\tecxt\text{d}t}{\tecxt\text{d}\lambda}}&\propto \kla\left( {1-\frac{R_S}{r}} \right)\frac{\tecxt\text{d}t}{\tecxt\text{d}\lambda}=\tecxt\text{const}=: \epsilon\\
\frac{\tecxt\text{d}}{\tecxt\text{d}\lambda}\frac{\partial L}{\partial \frac{\tecxt\text{d}\phi}{\tecxt\text{d}\lambda}}&=\frac{\partial L}{\partial \phi}=0&\Rightarrow &&\frac{\partial L}{\partial \frac{\tecxt\text{d}\phi}{\tecxt\text{d}\lambda}}&\propto r^2\frac{\tecxt\text{d}\phi}{\tecxt\text{d}\lambda}=\tecxt\text{const}=:h
\end{align*}
Es gilt außerdem $L=\frac{\tecxt\text{d}s^2}{\tecxt\text{d}\lambda^2}=0$. Einsetzen der Konstanten und Umstellen ergibt dann das effektive Potential:
\begin{align*}
0&=\kla\left( {1-\frac{R_S}{r}} \right)^{-1}\epsilon^2-\kla\left( {1-\frac{R_S}{r}} \right)^{-1}\kla\left( {\frac{\tecxt\text{d}r}{\tecxt\text{d}s}} \right)^2-\frac{h^2}{r^2}\\
\Rightarrow\ \ \ \epsilon^2&=\kla\left( {\frac{\tecxt\text{d}r}{\tecxt\text{d}\lambda}} \right)^2+\kla\left( {1-\frac{R_S}{r}} \right)\frac{h^2}{r^2}=:\kla\left( {\frac{\tecxt\text{d}r}{\tecxt\text{d}\lambda}} \right)^2+V_\tecxt\text{eff}(r)
\end{align*}
Um die Bahnkurve $r=r(\phi)$ zu bestimmen, nutzt man die Kettenregel aus und substituiert dann $u=:\frac{1}{r}$:
\begin{align*}
\epsilon^2&=\kla\left( {\frac{\tecxt\text{d}r}{\tecxt\text{d}\phi}\frac{\tecxt\text{d}\phi}{\tecxt\text{d}\lambda}} \right)^2+V_\tecxt\text{eff}=\kla\left( {\frac{\tecxt\text{d}r}{\tecxt\text{d}\phi}} \right)^2\frac{h^2}{r^4}+\kla\left( {1-\frac{R_S}{r}} \right)\frac{h^2}{r^2}\\
\Rightarrow\ \ \ \epsilon^2&=\kla\left( {-\frac{1}{u^2}\frac{\tecxt\text{d}u}{\tecxt\text{d}\phi}} \right)^2h^2u^4+(1-uR_S)h^2u^2=h^2(u')^2+(1-u)h^2u^2
\end{align*}
Ableiten nach $\phi$ und dividieren durch $2h^2u'$ ergibt dann
\begin{align*}
u''+u&=\frac{3}{2}R_S u^2\numberthis\label{u double prime}
\end{align*}
Es soll nun eine Lösung in erster Ordnung von $R_S$ gefunden werden. Schreibe dazu $u=u_0+R_S u_1+\mathcal{O}(R_S^2)$ mit $u_0''+u_0=0$. Dann findet man zunächst die homogene Lösung $u_0=a\cos \phi$. Einsetzen von $u_0$ in (\ref{u double prime}) liefert
\begin{align*}
R_S(u_1''+u_1)&=\frac{3}{2}R_S(u_0+R_Su_1)^2+\mathcal{O}(R_S^2)\\
\Rightarrow\ \ \ u_1''+u_1&\approx\frac{3}{2}u_0^2=\frac{3a^2}{4}\kla\left( {1+\cos 2\phi} \right)\numberthis\label{u1 double prime}
\end{align*}
Der Ansatz $u_1=\frac{3a^2}{4}+n\cos 2\phi $ führt zu einer Lösung für $n$. Einsetzen in (\ref{u1 double prime}) und dividieren durch $\cos 2\phi$ ergibt:
\begin{align*}
-4n+n&=\frac{3a^2}{4}&\Rightarrow &&n&=-\frac{a^2}{4}
\end{align*}
Die vollständige Lösung lautet dann
\begin{align*}
u(\phi)&= a\cos \phi+\frac{3a^2R_S}{4}-\frac{a^2R_S}{4}\cos 2\phi +\mathcal{O}(R_S^2)=\\
&= a^2R_S+a\cos \phi-\frac{a^2R_S}{2}\cos^2\phi+\mathcal{O}(R_S^2)\\
\Rightarrow\ \ \ r(\phi)&\approx\frac{1/a}{aR_S+\cos \phi-\frac{aR_S}{2}\cos^2\phi},
\end{align*}
wobei $\cos 2\phi=2\cos^2\phi -1$ verwendet wurde. Einsetzen von $\phi=0$ in die Bahnkurve liefert dann einen Zusammenhang zwischen dem minimalen Abstand $r(0)=r_0$ und $a$:
\begin{align*}
r_0&=\frac{1/a}{\frac{aR_S}{2}+1}&\Rightarrow &&\frac{1}{r_0}&=\frac{a^2R_S}{2}+a\\
\Rightarrow\ \ \ 0&=a^2+\frac{2a}{R_S}-\frac{2}{R_Sr_0}&\Rightarrow &&a&=\frac{1}{R_S}\kla\left( {\sqrt{1+\frac{2R_S}{r_0}}-1} \right)=\frac{1}{r_0}-\frac{R_S}{4r_0^2}+\mathcal{O}(R_S^2)
\end{align*} 
Mit $b:=\frac{R_S}{r_0}<1$ kann dieses Ergebnis auch als $aR_S=\sqrt{1+2b}-1= b+\mathcal{O}(R_S^2)$ in die Bahnkurve eingesetzt werden:
\begin{align*}
r(\phi)&\approx\frac{R_S/b}{b+\cos\phi -\frac{b}{2}\cos^2\phi}
\end{align*}
Wir wollen nun asymptotische Lösungen mit $r(\phi_{\tecxt\text{in},\tecxt\text{out}})=\infty$ finden. Mit $x:=\cos\phi$ lauten die Nullstellen des Nenners von $r(\phi)$ dann
\begin{align*}
0&\overset{!}{=} \frac{b}{2} x^2-b-x&\Rightarrow && x=\frac{1}{b}\kla\left( {1-\sqrt{1+2b^2}} \right)=b+\mathcal{O}(R_S^2)
\end{align*}
Das positive Vorzeichen liefert $x>1$ und ist somit keine Lösung, da $x=\cos\phi\overset{!}{<}1$. Man findet also $\phi_\tecxt\text{out}=-\phi_\tecxt\text{in}=\arccos (-b)+\mathcal{O}(R_S^2)$. Da die Formel ohnehin nur für kleine $b$ gilt, kann hier noch eine lineare Näherung gemacht werden. Mit $\arccos x=\frac{\pi}{2}+x+\mathcal{O}(x^2)$ ergibt sich der Ablenkwinkel dann zu
\begin{align*}
\Delta\phi &:=\phi_\tecxt\text{in}-\phi_\tecxt\text{out}-\pi=2\arccos b-\pi =2b+\mathcal{O}(R_S^2)\approx\frac{2GM}{r_0c^2}
\end{align*}
Für $r_0=7\cdot 10^8\,$m, $M=2\cdot 10^{30}\,$kg, $G=6.67\cdot 10^{-11}\,\frac{\tecxt\text{m}^3}{\tecxt\text{kg}\cdot\tecxt\text{s}^2}$ und $c=3\cdot 10^8\,\frac{\tecxt\text{m}}{\tecxt\text{s}}$ ergibt sich $\Delta\phi\approx 1.75''$.
\section*{Aufgabe 17: Lense-Thirring-Effekt}
In der Vorlesung wurde die Näherung $-\frac{\epsilon}{2}\square h_{\mu\nu}^\tecxt\text{E}=\kappa T_{\mu\nu}+\mathcal{O}(\epsilon^2)$ in der Eichung $\partial^\mu h_{\mu\nu}^\tecxt\text{E}=0$ hergeleitet. Für zeitunabhängige Quellen folgt daraus
\begin{align*}
\frac{\epsilon}{2}\nabla^2h_{0i}^\tecxt\text{E}&=\kappa T_{0i}=-\kappa\nabla^2A_i
\end{align*}
Gegeben ist ein Vektorpotential $\vec{A}=\frac{1}{2}\vec{B}\times\vec{r}$ mit $\vec{B}=\frac{1}{3}R^2\rho_0\vec{\omega}$. Schreibe die Komponenten $(h_{0i}^\tecxt\text{E})_i=:\vec{h}_0^\tecxt\text{E}$ als Vektor:
\begin{align*}
\epsilon \vec{h}_{0}^\tecxt\text{E}&=-2\kappa \vec{A}=\frac{1}{3}\kappa R^2\rho_0\vec{r}\times\vec{\omega}
\end{align*}
Dabei definiert man $h_{\mu\nu}^\tecxt\text{E}:=h_{\mu\nu}-\frac{1}{2}\eta_{\mu\nu}h$, wobei $\eta$ die Minkowski-Metrik ist. Daraus folgt aber direkt $h_{0i}^\tecxt\text{E}=h_{0i}$, da $\eta_{0i}=0$. Setze $\vec{\omega}=\omega\vec{e}_z$. Dann gilt $\epsilon h_{0x}=\frac{1}{3}\kappa R^2\rho_0\omega y$, $\epsilon h_{0y}=-\frac{1}{3}\kappa R^2\rho_0\omega x$ und $\epsilon h_{0z}=0$. Die Christoffel-Symbole lauten dann
\begin{align*}
\Gamma^j_{0i}&=\frac{\epsilon}{2}\eta^{j\sigma}\kla\left( {h_{\sigma 0,i}+\cancelto{0}{h_{\sigma i,0}}-h_{0i,\sigma}} \right)+\mathcal{O}(\epsilon^2)=-\frac{\epsilon}{2}\kla\left( {h_{0j,i}-h_{0i,j}} \right)+\mathcal{O}(\epsilon^2)
\end{align*}
Aufgrund der Symmetrie sind nur Terme mit $i\neq j$ ungleich Null. Letztlich erhält man $\Gamma^x_{0y}=\Gamma^x_{y0}=-\frac{1}{3}\kappa R^2\rho_0\omega$ und $\Gamma^y_{0x}=\Gamma^y_{x0}=\frac{1}{3}\kappa R^2\rho_0\omega$. Aus Aufgabe 14 wissen wir aber, dass $\Gamma^x_{y0}\equiv -\Omega$ und $\Gamma^y_{x0}\equiv\Omega$, woraus $\Omega=\frac{1}{3}\kappa R^2\rho_0\omega$ folgt. Für $\kappa =\frac{8\pi G}{c^2}$ und $\omega=\frac{2\pi}{T}$ mit $T=25\,$Tagen erhält man dann $\Omega=\frac{8\pi G}{3c^2}\frac{MR^2}{\frac{4}{3}\pi R^3}\frac{2\pi}{T}=\frac{2\pi GM}{RTc^2}=8.5\cdot 10^{-8}\,\frac{2\pi}{\tecxt\text{Tag}}$.
\end{document}